\documentclass[11pt]{article}

% --------------------------------------------------
% PACKAGES
% --------------------------------------------------
\usepackage[a4paper,margin=1in]{geometry}
\usepackage{amsmath,amssymb,amsfonts,amsthm}
\usepackage{mathrsfs}
\usepackage{mathtools}
\usepackage{enumitem}
\usepackage{hyperref}
\usepackage{color}
\usepackage{tikz-cd}
\usepackage{fancyhdr}

% --------------------------------------------------
% HYPERREF SETUP
% --------------------------------------------------
\hypersetup{
    colorlinks=true,
    linkcolor=blue,
    citecolor=blue,
    urlcolor=blue
}

% --------------------------------------------------
% THEOREM ENVIRONMENTS
% --------------------------------------------------
\newtheorem{theorem}{Theorem}[section]
\newtheorem{proposition}[theorem]{Proposition}
\newtheorem{lemma}[theorem]{Lemma}
\newtheorem*{theoremA}{\textbf{Theorem A}}
\newtheorem*{theoremB}{\textbf{Theorem B}}
\newtheorem*{corollaryA}{\textbf{Corollary A}}
\newtheorem*{corollaryB}{\textbf{Corollary B}}
\newtheorem*{corollaryc}{\textbf{Corollary C}}
\newtheorem{corollary}[theorem]{Corollary}
\newtheorem{conjecture}[theorem]{Conjecture}

\theoremstyle{definition}
\newtheorem{definition}[theorem]{Definition}
\newtheorem{example}[theorem]{Example}

\theoremstyle{remark}
\newtheorem{remark}[theorem]{Remark}

\newtheorem*{corollaryC}{\textbf{Corollary C}}

% --------------------------------------------------
% CUSTOM COMMANDS
% --------------------------------------------------
\newcommand{\Db}{\mathsf{D}^b}
\newcommand{\Dsg}{\mathsf{D}_{\mathrm{sg}}}
\newcommand{\modu}{\text{-mod}}
\newcommand{\gmod}{\text{-gmod}}
\newcommand{\proj}{\text{-proj}}
\newcommand{\inj}{\text{-inj}}

\newcommand{\Hom}{\operatorname{Hom}}
\newcommand{\Ext}{\operatorname{Ext}}

% --------------------------------------------------
% RUNNING HEADERS
% --------------------------------------------------
\setlength{\headheight}{14pt}
\pagestyle{fancy}
\fancyhf{}
\fancyhead[L]{\small\textsc{Koszul duality for absolutely Koszul algebras}}
\fancyhead[R]{\small\textsc{A. M. Bouhada}}
\fancyfoot[C]{\thepage}

% --------------------------------------------------
% TITLE INFORMATION
% --------------------------------------------------
\title{Koszul duality for finite-dimensional absolutely Koszul algebras}

\author{
A. M. Bouhada 
\date{}
}


\begin{document}

\maketitle
\pagestyle{plain}

% --------------------------------------------------
\begin{abstract}
Let \(\Lambda\) be a finite-dimensional Koszul algebra. Using a description of the linear part of minimal projective resolutions in terms of Koszul duality, we prove that \(\Lambda\) is absolutely Koszul if and only if its Koszul dual \(\Lambda^{!}\) is graded left co-coherent. As a consequence, every finite-dimensional graded module over a finite-dimensional absolutely Koszul algebra has a rational Poincar\'e series.
We further show that an absolutely Koszul algebra of finite global linearity defect has finite graded finitistic dimension.

Finally, we refine Koszul duality for finite-dimensional absolutely Koszul algebras. As a further application, we answer a question of Green et al.~\cite{12} by means of the Koszul dual algebra.
\end{abstract}

% --------------------------------------------------
\begin{center}
\section{Introduction}
\end{center}

Absolutely Koszul algebras form an important class of Koszul algebras. They were introduced by Iyengar and R\"omer in the commutative setting~\cite{15}; see also~\cite{14}. They are closely related to the study of the linear part of free resolutions over exterior algebras developed by Eisenbud, Fl{\o}ystad, and Schreyer~\cite{11}. The corresponding homological condition had already appeared in the noncommutative setting, notably in the work of Mart\'inez-Villa and Zacharia~\cite{19}, although the terminology \emph{absolutely Koszul} was not used there. We adopt this terminology throughout, while keeping in mind that absolutely Koszul algebras in the commutative setting satisfy additional properties that need not persist in the noncommutative case.

Despite the extensive study of absolute Koszulity in commutative algebra, its noncommutative counterpart remains less well understood, particularly from the viewpoint of the representation theory of finite-dimensional algebras. In our previous work~\cite{5}, we proved that quadratic monomial algebras are absolutely Koszul. More precisely, if \(\Lambda\) is a quadratic monomial algebra, then its global linearity defect is at most one. This contrasts with the commutative case, where quadratic monomial algebras need not be absolutely Koszul; see, for instance,~\cite{10}.

Our interest in these algebras is primarily motivated by the derived and singular Koszul dualities established in~\cite{4}. For a finite-dimensional Koszul algebra \(\Lambda\), we constructed triangulated equivalences
\[
\begin{tikzcd}[column sep=huge]
\mathsf{D}^{b}\!\bigl(
\Lambda^{!}\textup{-}\mathrm{Cop}^{\mathbb{Z}}
\bigr)
\arrow[r, bend left=35, "\mathfrak{F}"]
&
\mathsf{D}^{b}\!\bigl(
\Lambda\textup{-}\mathrm{gmod}
\bigr)
\arrow[l, bend left=35, "\mathfrak{G}"']
\end{tikzcd}
\]
and
\[
\begin{tikzcd}[column sep=huge]
\mathsf{D}^{b}\!\left(
\Lambda^{!}\textup{-}\mathrm{Cop}^{\mathbb{Z}}
\big/
\Lambda^{!}\textup{-}\mathrm{gmod}
\right)
\arrow[r, bend left=35, "\mathfrak{F}"]
&
\mathsf{D}_{\mathrm{sg}}\!\bigl(
\Lambda\textup{-}\mathrm{gmod}
\bigr)
\arrow[l, bend left=35, "\mathfrak{G}"']
\end{tikzcd}
\]
where \(\Lambda^{!}\textup{-}\mathrm{Cop}^{\mathbb{Z}}\) denotes the exact category of coperfect graded \(\Lambda^{!}\)-modules, namely, graded modules admitting finite injective coresolutions whose terms are finitely cogenerated injective modules.

These equivalences naturally lead to the problem of understanding the intrinsic structure of the category
\(\Lambda^{!}\textup{-}\mathrm{Cop}^{\mathbb{Z}}\).
This category is abelian precisely when the Koszul dual \(\Lambda^{!}\) is graded left co-coherent. In that case,
\(\Lambda^{!}\textup{-}\mathrm{Cop}^{\mathbb{Z}}\)
coincides with the category of finitely copresented graded \(\Lambda^{!}\)-modules, and one may form the abelian quotient
\(
\Lambda^{!}\textup{-}\mathrm{Cop}^{\mathbb{Z}}
\big/
\Lambda^{!}\textup{-}\mathrm{gmod}.
\)
This quotient is of particular interest in noncommutative geometry; see~\cite{22,23}. It plays a role analogous to that of the tails category in the construction of noncommutative projective schemes in the noetherian setting; see~\cite{2}. In~\cite{5}, we showed that quadratic monomial algebras satisfy the relevant coherence and co-coherence conditions.

These observations suggest a close connection between homological properties of \(\Lambda\) and co-coherence properties of its Koszul dual. We are therefore led to the following question.

\medskip

\noindent\textbf{Question.}
\emph{For which finite-dimensional Koszul algebras \(\Lambda\) is the category
\(\Lambda^{!}\textup{-}\mathrm{Cop}^{\mathbb{Z}}\) abelian?
Equivalently, when is \(\Lambda^{!}\) graded left co-coherent?
Can this condition be characterized intrinsically in terms of homological properties of \(\Lambda\), and in particular in terms of absolute Koszulity?}

\medskip

Related questions were studied by Yanagawa~\cite{26,27}, who used Koszul duality~\cite{3} to relate linearity defect and regularity over \(\Lambda\) to homological properties of modules over its Koszul dual. Recall that \(\Lambda\) is absolutely Koszul precisely when every finite-dimensional graded \(\Lambda\)-module has finite linearity defect. Yanagawa proved, in particular, that absolute Koszulity implies a coherence property for the Koszul dual. His convention for Koszul duality differs from ours by passage to the opposite algebra, whereas our convention agrees with that of~\cite{3}.

Conversely, Mart\'inez-Villa and Zacharia~\cite{19} proved absolute Koszulity under the stronger assumption that the Koszul dual is graded left noetherian. Thus, their result gives a converse to the implication obtained by Yanagawa, but under a substantially stronger finiteness hypothesis. The theorem below shows that noetherianity is not necessary: in our convention, graded left co-coherence is the appropriate condition.

Our approach is based on an explicit realization of graded derived Koszul duality. This allows us to identify absolute Koszulity directly with a finiteness property on the Koszul dual side and yields the following characterization.

\begin{theoremA}
Let \(\Lambda\) be a finite-dimensional Koszul algebra. Then \(\Lambda\) is absolutely Koszul if and only if its Koszul dual \(\Lambda^{!}\) is graded left co-coherent.
\end{theoremA}

\medskip

This result also sheds light on coherence phenomena for Koszul algebras of finite global dimension; compare with~\cite[Conjecture~1.3]{22}.

\medskip

Our second main result concerns the graded finitistic dimension. Using the linear part of minimal projective resolutions together with Koszul duality, we obtain the following finiteness criterion.

\begin{corollaryA}
Let \(\Lambda\) be a finite-dimensional absolutely Koszul algebra. If the global linearity defect of \(\Lambda\) is finite, then the graded finitistic dimension of \(\Lambda\) is finite.
\end{corollaryA}

\medskip

Finally, Theorem~A allows us to refine the graded and ungraded Koszul dualities established in~\cite{4}.

\begin{corollaryB}
Let \(\Lambda\) be a finite-dimensional absolutely Koszul algebra. Then the following triangulated equivalences hold.

\begin{enumerate}
\itemsep=0.4em

\item \textbf{Graded derived Koszul duality.}
\[
\mathsf{D}^{b}\!\bigl(
\Lambda^{!}\textup{-}\mathrm{Fp}^{\mathbb{Z}}
\bigr)
\xrightarrow{\ \sim\ }
\mathsf{D}^{b}\!\bigl(
\Lambda^{!}\textup{-}\mathrm{Fcp}^{\mathbb{Z}}
\bigr)
\xrightarrow{\ \sim\ }
\mathsf{D}^{b}\!\bigl(
\Lambda\textup{-}\mathrm{gmod}
\bigr).
\]

\item \textbf{Graded singular Koszul duality.}
\[
\mathsf{D}^{b}\!\left(
\Lambda^{!}\textup{-}\mathrm{Fcp}^{\mathbb{Z}}
\big/
\Lambda^{!}\textup{-}\mathrm{gmod}
\right)
\xrightarrow{\ \sim\ }
\mathsf{D}_{\mathrm{sg}}\!\bigl(
\Lambda\textup{-}\mathrm{gmod}
\bigr).
\]

\item \textbf{Graded BGG correspondence.}
If \(\Lambda\) is Iwanaga--Gorenstein, then
\[
\mathsf{D}^{b}\!\left(
\Lambda^{!}\textup{-}\mathrm{Fp}^{\mathbb{Z}}
\big/
\Lambda^{!}\textup{-}\mathrm{gmod}
\right)
\xrightarrow{\ \sim\ }
\mathsf{D}^{b}\!\left(
\Lambda^{!}\textup{-}\mathrm{Fcp}^{\mathbb{Z}}
\big/
\Lambda^{!}\textup{-}\mathrm{gmod}
\right)
\xrightarrow{\ \sim\ }
\Lambda\textup{-}\underline{\operatorname{Gproj}}^{\mathbb{Z}}.
\]

\item \textbf{Ungraded derived Koszul duality.}
\[
\mathrm{H}^{0}\!\left(
\operatorname{pretr}\!\left(
\mathsf{D}_{\mathrm{dg}}^{b}\!\bigl(
\Lambda^{!}\textup{-}\mathrm{Fcp}^{\mathbb{Z}}
\bigr)
\big/
\langle 1\rangle[1]
\right)
\right)
\xrightarrow{\ \sim\ }
\mathsf{D}^{b}\!\bigl(
\Lambda\textup{-}\mathrm{mod}
\bigr).
\]

\item \textbf{Ungraded singular Koszul duality.}
\[
\mathrm{H}^{0}\!\left(
\operatorname{pretr}\!\left(
\mathsf{D}_{\mathrm{dg}}^{b}\!\left(
\Lambda^{!}\textup{-}\mathrm{Fcp}^{\mathbb{Z}}
\big/
\Lambda^{!}\textup{-}\mathrm{gmod}
\right)
\big/
\langle 1\rangle[1]
\right)
\right)
\xrightarrow{\ \sim\ }
\mathsf{D}_{\mathrm{sg}}\!\bigl(
\Lambda\textup{-}\mathrm{mod}
\bigr).
\]

\item \textbf{Ungraded BGG correspondence.}
If \(\Lambda\) is Iwanaga--Gorenstein, then
\[
\mathrm{H}^{0}\!\left(
\operatorname{pretr}\!\left(
\mathsf{D}_{\mathrm{dg}}^{b}\!\left(
\Lambda^{!}\textup{-}\mathrm{Fcp}^{\mathbb{Z}}
\big/
\Lambda^{!}\textup{-}\mathrm{gmod}
\right)
\big/
\langle 1\rangle[1]
\right)
\right)
\xrightarrow{\ \sim\ }
\Lambda\textup{-}\underline{\operatorname{Gproj}}.
\]

\end{enumerate}

Here
\(\mathrm{H}^{0}\!\bigl(\operatorname{pretr}(-)\bigr)\)
denotes the homotopy category of the pretriangulated hull of the corresponding dg orbit category. For the original construction and proofs of these equivalences, we refer to~\cite{4}.
\section*{Acknowledgments}

Most of the ideas and results developed in this paper originated during the author's visit to the University of South Florida in the spring of 2020. The author would like to thank M. Elhamdadi for his warm hospitality.
\end{corollaryB}

\section{Preliminaries and Notation}
Throughout this paper, unless otherwise specified, $\Lambda$ denotes a finite-dimensional Koszul algebra of the form $kQ/I$, and $\Lambda^{!}$ denotes its Koszul dual. We regard both $\Lambda$ and $\Lambda^{!}$ as $k$-linear categories, or as $k$-linear graded categories when the grading is relevant. Accordingly, modules over $\Lambda$ and $\Lambda^{!}$ are viewed as $k$-linear functors.

\medskip

Since \(\Lambda\) is finite-dimensional, we work primarily with finite-dimensional modules. A \emph{left \(\Lambda\)-module} is a covariant \(k\)-linear functor
\[
M \colon \Lambda \longrightarrow k\textup{-mod},
\]
where \(k\textup{-mod}\) denotes the category of finite-dimensional \(k\)-vector spaces. Explicitly, such a functor assigns to each vertex \(x \in Q_0\) a vector space \(M(x)\), and to each arrow \(\alpha\colon x\to y\) a linear map
\[
M(\alpha)\colon M(x)\longrightarrow M(y),
\]
compatible with composition. Morphisms are given by natural transformations.

\medskip

We endow \(\Lambda\) with the grading induced by path length. A \emph{graded left \(\Lambda\)-module} is a covariant \(k\)-linear functor
\[
M \colon \Lambda \longrightarrow k\textup{-gmod},
\]
where \(k\textup{-gmod}\) denotes the category of finite-dimensional \(\mathbb{Z}\)-graded vector spaces. Thus
\[
M(x)=\bigoplus_{i\in\mathbb{Z}} N_i(x),
\]
and each morphism \(M(\alpha)\) is homogeneous of degree \(1\).

\medskip

Dually, the Koszul dual \(\Lambda^{!}\) is, in general, only locally finite-dimensional. Accordingly, we consider \emph{locally finite-dimensional graded \(\Lambda^{!}\)-modules}, that is, covariant \(k\)-linear functors
\[
M \colon \Lambda^{!} \longrightarrow k\textup{-GMod},
\]
where \(k\textup{-GMod}\) denotes the category of \(\mathbb{Z}\)-graded vector spaces such that each graded component is finite-dimensional.

\medskip

A graded module (over either \(\Lambda\) or \(\Lambda^{!}\)) is said to be
\begin{itemize}
\item \emph{left bounded} if it vanishes in sufficiently negative degrees,
\item \emph{right bounded} if it vanishes in sufficiently positive degrees,
\item \emph{bounded} if it is both left and right bounded.
\end{itemize}

\medskip

We denote by \(\Lambda\textup{-gmod}\) the category of finite-dimensional graded \(\Lambda\)-modules. For \(\Lambda^{!}\), we denote by \(\Lambda^{!}\textup{-GMod}\) the category of locally finite-dimensional graded modules, and by \(\Lambda^{!}\textup{-GMod}^{+}\) (resp.\ \(\Lambda^{!}\textup{-GMod}^{-}\)) the full subcategory of left bounded (resp.\ right bounded) modules.

\medskip

All these categories are naturally endowed with the grading shift functors
\[
\langle i\rangle \colon M \longmapsto M\langle i\rangle,
\qquad i \in \mathbb{Z},
\]
defined by reindexing the grading according to
\[
\bigl(M\langle i\rangle\bigr)_{n} = M_{n-i}.
\]

\medskip

For each vertex \(x\in Q_{0}\), the indecomposable graded projective \(\Lambda\)-module is given by
\[
P_{x}(y)=\Lambda(x,y),
\]
with grading induced by path length. The action of an arrow \(\alpha\colon a\to b\) is given by right multiplication,
\[
P_{x}(\alpha)(p)=p\alpha.
\]
We denote by \(\Lambda\textup{-proj}^{\mathbb{Z}}\) the category of finite-dimensional graded projective \(\Lambda\)-modules.

\medskip

Dually, for each vertex \(x\in Q_{0}\), the indecomposable graded injective \(\Lambda\)-module is given by
\[
I_{x}(y)=D\Lambda(y,x),
\qquad D=\operatorname{Hom}_{k}(-,k),
\]
with structure maps induced by precomposition,
\[
I_{x}(\alpha)(f)(p)=f(p\alpha).
\]
We denote by \(\Lambda\textup{-inj}^{\mathbb{Z}}\) the category of finite-dimensional graded injective \(\Lambda\)-modules.

\medskip

\medskip

Analogous constructions are carried out for the Koszul dual algebra \(\Lambda^{!}\), where one replaces finite-dimensional modules by locally finite-dimensional graded modules. In particular, we denote by
\( 
\Lambda^{!}\textup{-Proj}^{\mathbb{Z}}
\)
and \( 
\Lambda^{!}\textup{-Inj}^{\mathbb{Z}}
\)
the full subcategories of \(\Lambda^{!}\textup{-GMod}\) consisting of finite direct sums of grading shifts of indecomposable projective and injective modules, respectively.

\medskip

We write \(\Lambda^{!}\textup{-Fp}^{\mathbb{Z}}\) for the category of finitely presented graded \(\Lambda^{!}\)-modules, that is, those modules \(M\) fitting into an exact sequence
\[
P^{-1} \longrightarrow P^{0} \longrightarrow M \longrightarrow 0,
\]
with \(P^{-1}\) and \(P^{0}\) finite direct sums of objects in \(\Lambda^{!}\textup{-Proj}^{\mathbb{Z}}\).

Dually, we denote by \(\Lambda^{!}\textup{-Fcp}^{\mathbb{Z}}\) the category of finitely copresented graded \(\Lambda^{!}\)-modules, namely those modules \(N\) admitting an exact sequence
\[
0 \longrightarrow N \longrightarrow I^{0} \longrightarrow I^{1},
\]
where \(I^{0}\) and \(I^{1}\) belong to \(\Lambda^{!}\textup{-Inj}^{\mathbb{Z}}\).

\medskip

The algebra \(\Lambda^{!}\) is called \emph{left coherent} (respectively,
\emph{left co-coherent}) if the category
\(\Lambda^{!}\textup{-}\mathrm{Fp}^{\mathbb{Z}}\)
(respectively,
\(\Lambda^{!}\textup{-}\mathrm{Fcp}^{\mathbb{Z}}\))
is abelian.

Under the left coherence assumption, one may form the quotient category
\[
\Lambda^{!}\textup{-}\mathrm{Fp}^{\mathbb{Z}}
\big/
\Lambda^{!}\textup{-}\mathrm{gmod},
\]
usually referred to as the \emph{tails category}. It plays the role of the
category of coherent sheaves on a noncommutative projective space; see
~\cite{2,22,23}.

In the dual setting, if \(\Lambda^{!}\) is left co-coherent, one obtains the \emph{cotails category}
\[
\Lambda^{!}\textup{-Fcp}^{\mathbb{Z}} \big/ \Lambda^{!}\textup{-gmod}.
\]

\medskip

Even in the absence of coherence, the exact categories
\[
\Lambda^{!}\textup{-Cop}^{\mathbb{Z}} \big/ \Lambda^{!}\textup{-gmod}
\quad \text{and} \quad
\Lambda^{!}\textup{-Pe}^{\mathbb{Z}} \big/ \Lambda^{!}\textup{-gmod}
\]
remain well-defined and play a crucial role in the formulation of graded singular Koszul duality and variants of the Bernstein–Gelfand–Gelfand correspondence; see~\cite{4}.

\medskip

We denote by \(\Lambda^{!}\textup{-Pe}^{\mathbb{Z}}\) the category of \emph{perfect} graded \(\Lambda^{!}\)-modules, i.e., those admitting a finite projective resolution
\[
0 \longrightarrow P^{-n} \longrightarrow \cdots \longrightarrow P^{-1} \longrightarrow P^{0} \longrightarrow M \longrightarrow 0,
\]
with each \(P^{-i}\) a finite direct sum of objects from \(\Lambda^{!}\textup{-Proj}^{\mathbb{Z}}\).

Dually, the category \(\Lambda^{!}\textup{-Cop}^{\mathbb{Z}}\) consists of \emph{coperfect} graded modules, namely those modules admitting a finite injective coresolution
\[
0 \longrightarrow N \longrightarrow I^{0} \longrightarrow I^{1} \longrightarrow \cdots \longrightarrow I^{n} \longrightarrow 0,
\]
with each \(I^{i}\) in \(\Lambda^{!}\textup{-Inj}^{\mathbb{Z}}\).

\medskip


We recall that $\Lambda$ is Koszul if, for each vertex $x \in Q_{0}$, the simple module $S_{x}$ admits a linear graded projective resolution
\[
\cdots \longrightarrow P^{-2} \longrightarrow P^{-1} \longrightarrow P^{0} \longrightarrow S_{x} \longrightarrow 0,
\]
that is, each term $P^{-i}$ is generated in degree $-i$. Equivalently, one has
\[
P^{-i} \cong \bigoplus_{y \in Q_{0}} P_{y}\langle -i\rangle.
\]
Such a resolution is called \emph{linear}.

\medskip

More generally, one defines the notions of Koszul and coKoszul modules. A graded $\Lambda$-module (respectively, a graded $\Lambda^{!}$-module) $M$ is said to be \emph{Koszul} (respectively, \emph{coKoszul}) if it is generated in degree zero and admits a finite linear projective resolution (respectively, is cogenerated in degree zero and admits a finite colinear injective coresolution).

These conditions admit a homological characterization. A graded $\Lambda$-module $M$ is Koszul if
\[
\operatorname{Ext}^n_{\Lambda\textup{-gmod}}\bigl(M, S_x\langle -i\rangle\bigr)=0
\quad \text{for all } i \neq n \text{ and all } x \in Q_0,
\]
while a graded $\Lambda^{!}$-module $M$ is coKoszul if
\[
\operatorname{Ext}^n_{\Lambda^{!}\textup{-GMod}}\bigl(S^{!}_x\langle i\rangle, M\bigr)=0
\quad \text{for all } i \neq n \text{ and all } x \in Q_0.
\]

A graded module $M$ is said to be \emph{linear} (respectively, \emph{colinear}) if there exists an integer $i$ such that $M\langle -i\rangle$ is Koszul (respectively, coKoszul).

\medskip

The Koszul dual algebra $\Lambda^{!}$ may be described as follows. Let $V=(kQ)_2$ be the $k$-vector space spanned by all paths of length $2$ in $Q$, and set $I_2 = I \cap V$. Consider the opposite quiver $Q^{\mathrm{op}}$, and let $V^{\mathrm{op}} = (kQ^{\mathrm{op}})_2$. Choosing a basis of $V$ induces a dual basis of $V^{\mathrm{op}}$, and hence a nondegenerate bilinear pairing
\[
\langle -,- \rangle \colon V \times V^{\mathrm{op}} \longrightarrow k.
\]
Let $I_2^{\perp} \subset V^{\mathrm{op}}$ denote the orthogonal complement of $I_2$ with respect to this pairing. The quadratic dual algebra is then defined by
\[
\Lambda^{!} = kQ^{\mathrm{op}} \big/ \langle I_2^{\perp} \rangle.
\]

In general, even when $\Lambda$ is finite-dimensional, its Koszul dual $\Lambda^{!}$ is only locally finite-dimensional and need not be finite-dimensional. However, $\Lambda^{!}$ is finite-dimensional whenever $\Lambda$ has finite global dimension; see~\cite[Theorem~2.6]{4}. In fact, this condition is equivalent: $\Lambda$ has finite global dimension if and only if $\Lambda^{!}$ is finite-dimensional.

\medskip

We next recall the notion of the linear part of a minimal projective resolution, as well as the concepts of weakly Koszul modules and absolutely Koszul algebras.

Let \(M\) be a finitely presented graded \(\Lambda\)-module, and consider a minimal graded projective resolution
\[
\cdots \longrightarrow P^{-k} \longrightarrow \cdots \longrightarrow P^{-1} \longrightarrow P^{0} \longrightarrow M \longrightarrow 0.
\]
Since \(\Lambda\) is graded, each differential \(d^{i}\colon P^{i}\to P^{i+1}\) admits a decomposition into homogeneous components with respect to the internal grading,
\( 
d^{i}=\sum_{j\ge 1} d^{i}_{j},
\)
where \(d^{i}_{j}\) is homogeneous of degree \(j\).

\medskip

The \emph{linear part} of the resolution \(P^{\bullet}\), denoted \(\mathrm{Lin}^{\Lambda}(P^{\bullet})\), is the complex with the same underlying graded modules as \(P^{\bullet}\), but whose differentials are given by the degree-one components
\( 
d^{i}_{\mathrm{lin}} := d^{i}_{1}.
\)
Equivalently, \(\mathrm{Lin}^{\Lambda}(P^{\bullet})\) is obtained by discarding all higher-degree components of the differentials.

In the case where \(\Lambda\) is given by a quiver with relations, each component \(d^{i}_{j}\) arises from multiplication by paths of length \(j\). In particular, the linear part is induced by multiplication by arrows, and therefore \(\mathrm{Lin}^{\Lambda}(P^{\bullet})\) is again a complex of graded projective modules.

The notion of the linear part of a minimal projective resolution was introduced by Eisenbud, Fl{\o}ystad, and Schreyer in the context of exterior algebras; see~\cite{11}. Although originally formulated for the Koszul dual pair consisting of an exterior algebra and a polynomial algebra, the construction extends to arbitrary Koszul dual pairs. This perspective was further developed by Iwanaga, who generalized the theory to Koszul algebras and established several fundamental connections with Koszul duality; see~\cite{26}.

\medskip

The \emph{linearity defect} of a finitely presented graded \(\Lambda\)-module \(M\) is defined by
,
\[
\mathrm{ld}_{\Lambda}(M)
=
\sup\Bigl\{
i 
\;\Big|\;
H^{i}\bigl(\mathrm{Lin}^{\Lambda}(P^{\bullet})\bigr)\neq 0
\Bigr\}.
\]

\medskip

The \emph{global linearity defect} of \(\Lambda\) is given by
\[
\mathrm{gl\,ld}_{\Lambda}
=
\sup\Bigl\{
\mathrm{ld}_{\Lambda}(M)
\;\Big|\;
M \in \Lambda\textup{-gmod}
\Bigr\}.
\]

\medskip

A graded \(\Lambda\)-module \(M\) is said to be \emph{weakly Koszul} if the linear part \(\mathrm{Lin}^{\Lambda}(P^{\bullet})\) of its minimal graded projective resolution is exact; equivalently, if \(\mathrm{ld}_{\Lambda}(M)=0\).

\medskip

Following~\cite{19,26}, we adopt the terminology \emph{weakly Koszul} in order to distinguish this notion from other usages of the term \emph{Koszul module}, including that of~\cite{14}, where modules satisfying the above condition are referred to as Koszul. A Koszul duality for weakly Koszul modules was established in~\cite[Theorem~3.1]{16} and~\cite[Theorem~4.4]{26}, and continues to hold in the present setting.
\medskip

A Koszul algebra \(\Lambda\) is said to be \emph{absolutely Koszul} if every finite-dimensional \(\Lambda\)-module admits a weakly Koszul syzygy. It was shown in~\cite{5} that quadratic monomial algebras are absolutely Koszul and that their global linearity defect is at most one. Moreover, in the radical square zero case, every finite-dimensional module is weakly Koszul, and hence the global linearity defect vanishes. Exterior algebras also provide a fundamental class of absolutely Koszul algebras; see~\cite{11}.

\medskip

Dually, one defines the notion of a \emph{co-absolutely Koszul} algebra by requiring that every finite-dimensional graded $\Lambda^{!}$-module admits a weakly coKoszul cosyzygy.

\medskip 
We now recall the construction of the graded derived Koszul duality in the finite-dimensional setting; see~\cite{4} for a detailed account.

\medskip

A cochain complex of graded projective \(\Lambda\)-modules
\[
\cdots \longrightarrow P^{n-1} \longrightarrow P^{n} 
\longrightarrow P^{n+1} \longrightarrow \cdots
\]
is said to be \emph{linear} if every indecomposable direct summand of \(P^{n}\) is of the form \(P_{x}\langle n\rangle\) for some \(x \in Q_{0}\). Equivalently, each summand is generated in internal degree \(n\). We denote by
\( 
\mathcal{LC}\!\bigl(\Lambda\textup{-proj}^{\mathbb{Z}}\bigr)
\)
the category of linear complexes of graded projective modules.

By results of Martínez-Villa and Saorín~\cite[Theorem~2.4]{18}, and independently Mazorchuk, Ovsienko, and Stroppel~\cite[Theorem~12]{20}, there is an equivalence of abelian categories
\[
\Lambda^{!}\textup{-GMod}
\xrightarrow{\;\sim\;}
\mathcal{LC}\!\bigl(\Lambda\textup{-proj}^{\mathbb{Z}}\bigr),
\]
which restricts to
\[
\Lambda^{!}\textup{-gmod}
\xrightarrow{\;\sim\;}
\mathcal{LC}^{b}\!\bigl(\Lambda\textup{-proj}^{\mathbb{Z}}\bigr).
\]

This equivalence is realized by the \emph{Koszul functor}
\[
K \colon \Lambda^{!}\textup{-GMod} \longrightarrow
\mathcal{LC}\!\bigl(\Lambda\textup{-proj}^{\mathbb{Z}}\bigr),
\]
defined by
\[
K(M)^{n} = \bigoplus_{x \in Q_{0}} P_{x}\langle n \rangle \otimes_{k} M_{n}(x),
\]
with differential
\[
d^{n} = \sum_{\alpha\colon y \to x}
P_{\alpha} \otimes M(\alpha^{\operatorname{op}}),
\qquad P_{\alpha}(p)=p\alpha.
\]

\medskip

Let \(\Lambda^{!}\textup{-GMod}^{-,b}\) denote the full subcategory of modules \(M\) such that \(K(M)\) is right bounded and has bounded cohomology. Then \(K\) restricts to an equivalence of exact categories
\[
K \colon \Lambda^{!}\textup{-GMod}^{-,b}
\xrightarrow{\;\sim\;}
\mathcal{LC}^{-,b}\!\bigl(\Lambda\textup{-proj}^{\mathbb{Z}}\bigr).
\]

Although \(\Lambda^{!}\textup{-GMod}^{-,b}\) is not abelian in general, it admits a natural exact structure: it is closed under extensions and satisfies the two-out-of-three property for short exact sequences. In particular, it is stable under kernels of admissible epimorphisms and cokernels of admissible monomorphisms. Since \(\Lambda\) is finite-dimensional, one has
\[
\Lambda^{!}\textup{-GMod}^{-,b} = \Lambda^{!}\textup{-Cop}^{\mathbb{Z}},
\qquad \text{see~\cite{4}}.
\]

\medskip

We now pass to derived categories. Let
\[
M^{\bullet} \in \mathsf{C}^{b}\!\bigl(\Lambda^{!}\textup{-GMod}^{-,b}\bigr)
\]
be a bounded cochain complex. Applying \(K\) componentwise yields a double complex
\[
B^{q,p} := K(M^{p})^{q}
= \bigoplus_{x\in Q_{0}} P_{x}\langle q\rangle \otimes M^{p}_{q}(x).
\]

The horizontal differential is induced by \(d^{p}\), while the vertical differential is given by the \(\Lambda^{!}\)-action:
\[
d_{1}^{q,p} =
\bigoplus_{x} \mathrm{id} \otimes (d^{p})_{q,x},
\qquad
d_{2}^{q,p} =
\sum_{\alpha\colon x \to y}
P_{\alpha} \otimes M^{p}(\alpha^{\mathrm{op}}).
\]

The associated total complex is defined by
\[
\mathrm{Tot}(B)^{n} = \bigoplus_{p+q=n} B^{q,p},
\qquad
d_{\mathrm{Tot}} = d_{1} + (-1)^{p} d_{2},
\]
and satisfies \(d_{\mathrm{Tot}}^{2}=0\).

\medskip

This construction defines a functor
\[
\mathcal{F}\colon
\mathsf{C}^{b}\!\bigl(\Lambda^{!}\textup{-GMod}^{-,b}\bigr)
\longrightarrow
\mathsf{C}^{-,b}\!\bigl(\Lambda\textup{-proj}^{\mathbb{Z}}\bigr),
\]
which preserves homotopies. Passing to homotopy categories, we obtain
\[
\mathscr{F}\colon
\mathsf{K}^{b}\!\bigl(\Lambda^{!}\textup{-GMod}^{-,b}\bigr)
\longrightarrow
\mathsf{K}^{-,b}\!\bigl(\Lambda\textup{-proj}^{\mathbb{Z}}\bigr).
\]

Using the canonical equivalence
\[
\mathsf{K}^{-,b}\!\bigl(\Lambda\textup{-proj}^{\mathbb{Z}}\bigr)
\cong
\mathsf{D}^{b}\!\bigl(\Lambda\textup{-gmod}\bigr),
\]
together with the Acyclic Assembly Lemma, the functor \(\mathscr{F}\) descends to a triangulated functor
\[
\mathfrak{F}\colon
\mathsf{D}^{b}\!\bigl(\Lambda^{!}\textup{-Cop}^{\mathbb{Z}}\bigr)
\longrightarrow
\mathsf{D}^{b}\!\bigl(\Lambda\textup{-gmod}\bigr).
\]

\medskip

The quasi-inverse is constructed dually. For \(N \in \Lambda\textup{-gmod}\), define
\[
G(N)^{n}
=
\bigoplus_{x\in Q_{0}} I^{!}_{x}\langle -n\rangle \otimes N_{n}(x),
\qquad
I^{!}_{x} := \mathbb{D}\!\bigl(\Lambda^{!\,\mathrm{op}}(x,-)\bigr).
\]

The differential is induced by the \(\Lambda\)-module structure of \(N\). Extending to complexes yields a functor
\[
\mathcal{G}\colon
\mathsf{C}^{b}\!\bigl(\Lambda\textup{-gmod}\bigr)
\longrightarrow
\mathsf{C}^{b}\!\bigl(\Lambda^{!}\textup{-Inj}^{\mathbb{Z}}\bigr),
\]
which preserves homotopies and induces
\[
\mathfrak{G}\colon
\mathsf{D}^{b}\!\bigl(\Lambda\textup{-gmod}\bigr)
\longrightarrow
\mathsf{D}^{b}\!\bigl(\Lambda^{!}\textup{-Cop}^{\mathbb{Z}}\bigr).
\]

\medskip

The functors \(\mathfrak{F}\) and \(\mathfrak{G}\) are quasi-inverse triangulated equivalences. The functor \(\mathfrak{F}\) is referred to as the \emph{graded derived Koszul duality}; see~\cite{4}.
\begin{remark}
The graded derived Koszul duality constructed in this paper, together with its quasi-inverse, is closely related to the duality developed by Eisenbud, Fl{\o}ystad, and Schreyer in the context of exterior algebras; see~\cite{11}.
\end{remark}
\medskip 
The graded derived Koszul duality functor \(\mathfrak{F}\) and its quasi-inverse \(\mathfrak{G}\) may be viewed as restrictions of a more general equivalence between suitable triangulated subcategories of \(\mathsf{D}(\Lambda^{!}\textup{-GMod})\) and \(\mathsf{D}(\Lambda\textup{-GMod})\), namely
\[
\mathsf{D}^{\downarrow}(\Lambda^{!}\textup{-GMod})
\;\xrightarrow{\ \sim\ }\;
\mathsf{D}^{\uparrow}(\Lambda\textup{-GMod}).
\]
For a detailed construction, including precise definitions and proofs, we refer to~\cite[Theorem~5.7]{7}, \cite[Theorem~30]{20}, and \cite[Theorem~2.12.1]{3}; see also~\cite{8}. This equivalence will be used throughout the paper without further comment.
\medskip

The principal object of study in this paper is the triangulated equivalence
\[
\begin{tikzcd}[column sep=huge]
\mathsf{D}^b\bigl(\Lambda^{!}\textup{-Cop}^{\mathbb{Z}}\bigr)
\arrow[r, bend left=35, "\mathfrak{F}"]
&
\mathsf{D}^b\bigl(\Lambda\textup{-gmod}\bigr)
\arrow[l, bend left=35, "\mathfrak{G}"']
\end{tikzcd}
\]
Since the category \(\Lambda^{!}\textup{-Cop}^{\mathbb{Z}}\) is not abelian in general, we regard
\(\mathsf{D}^b\bigl(\Lambda^{!}\textup{-Cop}^{\mathbb{Z}}\bigr)\) as a full triangulated subcategory of
\(\mathsf{D}^{\downarrow}(\Lambda^{!}\textup{-GMod})\). In particular, we shall freely consider the cohomology objects of complexes in
\(\mathsf{D}^b\bigl(\Lambda^{!}\textup{-Cop}^{\mathbb{Z}}\bigr)\).
\begin{center}
\section{Main Results}
\end{center}
In this section, we prove the main results announced in the introduction. We begin by describing the linear part of minimal projective resolutions via graded derived Koszul duality and use this description to characterize finite-dimensional absolutely Koszul algebras in terms of their Koszul duals. We then derive a necessary and sufficient condition for the finiteness of the graded finitistic dimension of a finite-dimensional Koszul algebra. As an application, we show that absolutely Koszul algebras of finite global linearity defect have finite graded finitistic dimension. We next use these results to refine the Koszul dualities established in~\cite{4}. Finally, we answer a question raised by Green et al.~\cite{12} concerning when the category of modules with linear presentations coincides with the category of Koszul modules.

\medskip 

The following result is well known; see, for instance,~\cite{3,4,7,20}.

\begin{proposition}
The graded derived Koszul duality functor $\mathfrak{F}$ sends each indecomposable injective module $I^{!}_{x}$ to a minimal projective resolution of the corresponding simple module $S_{x}$, and sends simple modules to projective modules.

Dually, its quasi-inverse $\mathfrak{G}$ sends each indecomposable projective module $P_{x}$ to a minimal injective coresolution of the corresponding simple module $S^{!}_{x}$, and sends simple modules to injective modules.
\end{proposition}

Before recalling an explicit description of the linear part of a minimal projective resolution, we require the following auxiliary results concerning Koszul duality. The first lemma is taken from~\cite[Lemma~4.2]{4}. We note that the grading shift used in~\cite{4} is opposite to the convention adopted in the present paper.

\begin{lemma}
Let \(X^{\bullet} \in \mathsf{D}^{b}\!\bigl(\Lambda^{!}\textup{-Cop}^{\mathbb{Z}}\bigr)\). Then, for every \(i \in \mathbb{Z}\), there is a natural isomorphism
\[
\mathfrak{F}\bigl(X^{\bullet}\langle i\rangle\bigr)
\;\cong\;
\mathfrak{F}(X^{\bullet})\langle i\rangle [-i].
\]
\end{lemma}

\medskip

The following lemma gives an explicit description of the cohomology on the Koszul dual side in terms of morphisms into indecomposable injectives.

\begin{lemma}
Let
\[
X^{\bullet}\in
\mathsf{D}^{b}\!\bigl(
\Lambda^{!}\textup{-}\mathrm{Cop}^{\mathbb{Z}}
\bigr),
\]
and let \(I_x^{!}\) be the indecomposable graded injective
\(\Lambda^{!}\)-module corresponding to \(x\in Q_0\). Then, for all
\(m,i\in\mathbb Z\), there are natural isomorphisms of \(k\)-vector spaces
\[
\operatorname{Hom}_{\mathsf{D}^{b}
(\Lambda^{!}\textup{-}\mathrm{Cop}^{\mathbb{Z}})}
\bigl(
X^{\bullet},
I_x^{!}\langle m\rangle[-i]
\bigr)
\cong
\operatorname{Hom}_{\Lambda^{!}\textup{-}\mathrm{GMod}}
\bigl(
H^{i}(X^{\bullet}),
I_x^{!}\langle m\rangle
\bigr)
\cong
D\!\left(H^{i}(X^{\bullet})_{m}(x)\right).
\]
\end{lemma}

\begin{proof}
Set \(I=I_x^{!}\langle m\rangle\). Since \(I\) is injective, the stalk
complex \(I[-i]\), concentrated in cohomological degree \(i\), is
\(K\)-injective. Hence
\[
\operatorname{Hom}_{\mathsf{D}^{b}
(\Lambda^{!}\textup{-}\mathrm{Cop}^{\mathbb{Z}})}
\bigl(X^\bullet,I[-i]\bigr)
\cong
\operatorname{Hom}_{\mathsf{K}^{b}
(\Lambda^{!}\textup{-}\mathrm{Cop}^{\mathbb{Z}})}
\bigl(X^\bullet,I[-i]\bigr).
\]

A chain map \(X^\bullet\to I[-i]\) is determined by a morphism
\(f:X^i\to I\) satisfying
\[
f\circ d_X^{i-1}=0.
\]
Thus \(f\) factors uniquely through
\(X^i/\operatorname{Im}d_X^{i-1}\), and therefore
\[
Z^0\operatorname{Hom}^{\bullet}(X^\bullet,I[-i])
\cong
\operatorname{Hom}\!\left(
X^i/\operatorname{Im}d_X^{i-1},I
\right).
\]

A chain map is null-homotopic precisely when the corresponding morphism
factors through \(d_X^i:X^i\to X^{i+1}\). Since \(I\) is injective, every
morphism
\[
\operatorname{Im}d_X^i\longrightarrow I
\]
extends to \(X^{i+1}\). Hence the null-homotopic maps identify with the
image of
\[
\operatorname{Hom}(\operatorname{Im}d_X^i,I)
\longrightarrow
\operatorname{Hom}\!\left(
X^i/\operatorname{Im}d_X^{i-1},I
\right).
\]

Now the exact sequence
\[
0\longrightarrow H^i(X^\bullet)
\longrightarrow
X^i/\operatorname{Im}d_X^{i-1}
\longrightarrow
\operatorname{Im}d_X^i
\longrightarrow 0
\]
yields, by injectivity of \(I\),
\[
0\longrightarrow
\operatorname{Hom}(\operatorname{Im}d_X^i,I)
\longrightarrow
\operatorname{Hom}\!\left(
X^i/\operatorname{Im}d_X^{i-1},I
\right)
\longrightarrow
\operatorname{Hom}(H^i(X^\bullet),I)
\longrightarrow0.
\]
Consequently,
\[
\operatorname{Hom}_{\mathsf{D}^{b}
(\Lambda^{!}\textup{-}\mathrm{Cop}^{\mathbb{Z}})}
\bigl(X^\bullet,I_x^{!}\langle m\rangle[-i]\bigr)
\cong
\operatorname{Hom}_{\Lambda^{!}\textup{-}\mathrm{GMod}}
\bigl(H^i(X^\bullet),I_x^{!}\langle m\rangle\bigr).
\]

Finally, graded co-Yoneda and the convention
\((M\langle m\rangle)_n=M_{n-m}\) give
\[
\operatorname{Hom}_{\Lambda^{!}\textup{-}\mathrm{GMod}}
\bigl(H^i(X^\bullet),I_x^{!}\langle m\rangle\bigr)
\cong
D\!\left(H^i(X^\bullet)_m(x)\right).
\]
This proves the result.
\end{proof}

The following proposition exhibits a strong relationship between the minimal projective resolution of a module \(N\) and the cohomology of the associated complex in 
\(\mathsf{D}^{b}\!\bigl(\Lambda^{!}\textup{-Cop}^{\mathbb{Z}}\bigr)\).

\begin{proposition}
Let \(N\in\Lambda\textup{-}\mathrm{gmod}\), and let
\[
X^{\bullet}\in
\mathsf{D}^{b}\!\bigl(
\Lambda^{!}\textup{-}\mathrm{Cop}^{\mathbb{Z}}
\bigr)
\]
be such that
\[
\mathfrak{F}(X^{\bullet})\cong N
\qquad\text{in}\qquad
\mathsf{D}^{b}\!\bigl(
\Lambda\textup{-}\mathrm{gmod}
\bigr).
\]
Then, for every \(x\in Q_0\), every \(m\in\mathbb Z\), and every
\(i\geq0\), there is a natural isomorphism of \(k\)-vector spaces
\[
\operatorname{Ext}^{i}_{\Lambda\textup{-}\mathrm{gmod}}
\bigl(N,S_x\langle m\rangle\bigr)
\cong
D\!\left(
H^{-m-i}(X^{\bullet})_m(x)
\right).
\]
\end{proposition}
 \begin{proof}
By the realization of Ext in the bounded derived category,
\[
\operatorname{Ext}^{i}_{\Lambda\textup{-}\mathrm{gmod}}
\bigl(N,S_x\langle m\rangle\bigr)
\cong
\operatorname{Hom}_{\mathsf{D}^{b}
(\Lambda\textup{-}\mathrm{gmod})}
\bigl(N,S_x\langle m\rangle[i]\bigr).
\]
Since
\(\mathfrak{F}(X^{\bullet})\cong N\), we obtain
\[
\operatorname{Ext}^{i}_{\Lambda\textup{-}\mathrm{gmod}}
\bigl(N,S_x\langle m\rangle\bigr)
\cong
\operatorname{Hom}_{\mathsf{D}^{b}
(\Lambda\textup{-}\mathrm{gmod})}
\bigl(
\mathfrak{F}(X^{\bullet}),
S_x\langle m\rangle[i]
\bigr).
\]

By Proposition~3.1 and the compatibility of \(\mathfrak{F}\) with grading
shifts,
\[
\mathfrak{F}\bigl(I_x^{!}\langle m\rangle\bigr)
\cong
S_x\langle m\rangle[-m].
\]
Therefore
\[
S_x\langle m\rangle[i]
\cong
\mathfrak{F}\bigl(I_x^{!}\langle m\rangle\bigr)[i+m].
\]
Since \(\mathfrak{F}\) is fully faithful, it follows that
\[
\begin{aligned}
\operatorname{Hom}_{\mathsf{D}^{b}
(\Lambda\textup{-}\mathrm{gmod})}
\bigl(
\mathfrak{F}(X^{\bullet}),
S_x\langle m\rangle[i]
\bigr)
&\cong
\operatorname{Hom}_{\mathsf{D}^{b}
(\Lambda^{!}\textup{-}\mathrm{Cop}^{\mathbb{Z}})}
\bigl(
X^{\bullet},
I_x^{!}\langle m\rangle[i+m]
\bigr).
\end{aligned}
\]

Applying the preceding lemma with
\[
j=-(i+m),
\]
the latter space is naturally isomorphic to
\[
D\!\left(
H^{-(i+m)}(X^{\bullet})_m(x)
\right).
\]
Hence
\[
\operatorname{Ext}^{i}_{\Lambda\textup{-}\mathrm{gmod}}
\bigl(N,S_x\langle m\rangle\bigr)
\cong
D\!\left(
H^{-m-i}(X^{\bullet})_m(x)
\right),
\]
as claimed.
\end{proof}
The following result clarifies the relationship established in Proposition~3.4. It originates in the work of Eisenbud, Fl{\o}ystad, and Schreyer in the context of exterior algebras; see~\cite[Corollary~3.6]{11}, and extends naturally to finite-dimensional Koszul algebras. Yanagawa~\cite[Theorem~3.9]{26} further generalized this result to arbitrary pairs \((\Lambda,\Lambda^{!})\). It shows that the linear part of a minimal projective resolution of a graded \(\Lambda\)-module is completely determined by cohomological data over the Koszul dual.

\begin{proposition}
Let \(N\) be a finite-dimensional graded \(\Lambda\)-module, and let
\[
P^{\bullet} \longrightarrow N
\]
be its minimal graded projective resolution. Denote by \(\mathrm{Lin}^{\Lambda}(P^{\bullet})\) its linear part.

Then there exists a bounded colinear complex of injective \(\Lambda^{!}\)-modules
\[
I^{\bullet} \in \mathsf{D}^{b}\!\bigl(\Lambda^{!}\textup{-Cop}^{\mathbb{Z}}\bigr)
\]
such that
\[
\mathfrak{F}\!\left(\bigoplus_{i \in \mathbb{Z}} H^{i}(I^{\bullet})[-i]\right)
\;\cong\;
\mathrm{Lin}^{\Lambda}(P^{\bullet})
\]
in \(\mathsf{D}^{b}\!\bigl(\Lambda\textup{-gmod}\bigr)\).
\end{proposition}
\medskip

We are now in a position to establish the first main result of this paper. It refines ~\cite[Proposition~4.8]{27} and removes the noetherianity assumption imposed in~\cite[Theorem~4.5]{19}.

\begin{theorem}
A finite-dimensional Koszul algebra \(\Lambda\) is absolutely Koszul if and only if its Koszul dual \(\Lambda^{!}\) is graded left co-coherent.
\end{theorem}

\begin{proof}
Assume first that \(\Lambda\) is absolutely Koszul. By definition, every finite-dimensional graded \(\Lambda\)-module has finite linearity defect. Equivalently, every finite-dimensional graded module \(M\) has a weakly Koszul syzygy, or, what is the same, the linear part of its minimal projective resolution has bounded cohomology. By Proposition~3.5, it follows that
\[
\mathfrak{F}\!\left(\bigoplus_{i \in \mathbb{Z}} H^{i}(I^{\bullet})[-i]\right)
\;\cong\;
\mathrm{Lin}^{\Lambda}(P^{\bullet})
\]
has bounded cohomology. In particular, each object \(\mathfrak{F}(H^{i}(I^{\bullet})[-i])\) has bounded cohomology. By~\cite[Proposition~3.4]{4}, each \(H^{i}(I^{\bullet})\) is coperfect.

Now let \(M\) be a finitely copresented module. There exists an injective copresentation
\[
0 \longrightarrow M \longrightarrow J^{0} \longrightarrow J^{1}.
\]
Consider the complex
\[
J^{\bullet}\colon \cdots \longrightarrow 0 \longrightarrow J^{0} \longrightarrow J^{1} \longrightarrow 0 \longrightarrow \cdots.
\]
Clearly, \(\mathfrak{F}(J^{\bullet})\) has bounded cohomology. Since \(\mathfrak{F}(J^\bullet)\) is right bounded and has bounded cohomology, there exists \(n\in\mathbb Z\) such that
\[
H^i(\mathfrak{F}(J^\bullet))=0 \qquad \text{for all } i<n.
\]
Consider the stupid truncation triangle
\[
\sigma^{\le n}\mathfrak{F}(J^\bullet)
\longrightarrow
\mathfrak{F}(J^\bullet)
\longrightarrow
\sigma^{\ge n+1}\mathfrak{F}(J^\bullet)
\longrightarrow
\sigma^{\le n}\mathfrak{F}(J^\bullet)[1].
\]
The complex \(\sigma^{\le n}\mathfrak{F}(J^\bullet)\) is quasi-isomorphic to a shift of a module, since its cohomology vanishes in all degrees except possibly in degree \(n\). More precisely,
\[
\sigma^{\le n}\mathfrak{F}(J^\bullet)
\;\cong\;
H^{n}\!\bigl(\mathfrak{F}(J^\bullet)\bigr)[-n]
\]
in \(\mathsf{D}^{b}(\Lambda\textup{-gmod})\).

Since \(\Lambda\) is absolutely Koszul, the linear part of the minimal projective resolution of the module \(H^{n}(\mathfrak{F}(J^\bullet))\) has bounded cohomology. It follows that the cohomology objects of the bounded complex
\[
\mathfrak{G}\!\bigl(H^{n}(\mathfrak{F}(J^\bullet))[-n]\bigr)
\]
are coperfect graded modules.

Now apply \(\mathfrak{G}\) to the above exact triangle. We obtain an exact triangle in \(\mathsf{D}^{b}(\Lambda^{!}\textup{-Cop}^{\mathbb{Z}})\) of the form
\[
X^\bullet \longrightarrow J^\bullet \longrightarrow Y^\bullet \longrightarrow X^\bullet[1],
\]
such that \(Y^\bullet\) is a bounded complex with finite-dimensional cohomology and all the cohomology objects of \(X^\bullet\) are coperfect. Passing to the long exact sequence of cohomology and since \(\Lambda^{!}\textup{-Cop}^{\mathbb{Z}}\) satisfies the two-out-of-three property, it follows that
\[
H^{0}(J^\bullet)\cong M
\]
is a coperfect graded module. It follows that 
\(\Lambda^{!}\textup{-Cop}^{\mathbb{Z}}\) is abelian and therefore \(\Lambda^{!}\) is graded left co-coherent.
Conversely, assume that \(\Lambda^{!}\) is graded left co-coherent, so that the category \(\Lambda^{!}\textup{-Cop}^{\mathbb{Z}}\) is abelian. Let \(N\) be a finite-dimensional graded \(\Lambda\)-module. There exists a bounded colinear complex \(I^\bullet\) such that
\[
\mathfrak{F}(I^\bullet)\cong N.
\]
Since \(\Lambda^{!}\textup{-Cop}^{\mathbb{Z}}\) is abelian, all cohomology objects \(H^{i}(I^\bullet)\) are coperfect. Consequently,
\[
\mathfrak{F}(H^{i}(I^\bullet)[-i])
\]
has bounded cohomology for every \(i\). By Proposition~3.5, it follows that the linear part of the minimal projective resolution of \(N\) has bounded cohomology. Hence \(N\) has finite linearity defect. Therefore \(\Lambda\) is absolutely Koszul.
\end{proof}

We now record a first consequence of the preceding theorem. It concerns the
rationality of Poincar\'e series of finite-dimensional modules.

For a systematic investigation of Poincar\'e series of finite-dimensional
modules over Koszul algebras, Hilbert series, and their connection with Koszul
duality, we refer the reader to ~\cite{21}. 
\medskip

Recall that the Poincar\'e series of a
finite-dimensional \(\Lambda\)-module \(M\) is defined by
\[
P^{\Lambda}_{M}(t)
=
\sum_{n\geq 0}
\dim_{k}\operatorname{Ext}^{n}_{\Lambda}(M,\Lambda_{0})\,t^{n}.
\]
The following corollary generalizes~\cite[Theorem~4.7]{19}; the proof is
obtained by the same argument.

\begin{corollary}
Every finite-dimensional module over a finite-dimensional absolutely Koszul algebra has a rational Poincar\'e series.
\end{corollary}
\begin{remark}
The statement of~\cite[Theorem~5.8]{19}, concerning the stable components of the Auslander--Reiten quiver of a self-injective Koszul algebra of Loewy length greater than three, remains valid in the present setting.
\end{remark}

We now turn to a fundamental question in the homological algebra of finite-dimensional algebras. 

The classical finitistic dimension conjecture, formulated by \textsc{Bass} in 1960, asserts that for every finite-dimensional algebra \(\Lambda\), the finitistic dimension
\[
\mathrm{fin\,dim}\,\Lambda
=
\sup\{\mathrm{pd}_\Lambda M
\mid
M \in \Lambda\textup{-mod},\ \mathrm{pd}_\Lambda M < \infty\}
\]
is finite. Despite substantial progress and numerous partial results, this conjecture has remained open for more than six decades and is widely regarded as one of the central problems in the homological theory of finite-dimensional algebras.

To the best of our knowledge, the graded analogue of this invariant has not been systematically investigated. For a finite-dimensional graded algebra \(\Lambda\), one may define the graded finitistic dimension by
\[
\mathrm{grfin\,dim}\,\Lambda
=
\sup\{\mathrm{pd}_{\Lambda\textup{-gmod}} M
\mid
M\in \Lambda\textup{-gmod},\ \mathrm{pd}_{\Lambda\textup{-gmod}} M<\infty\}.
\]
In particular, the behavior of the graded finitistic dimension for finite-dimensional Koszul algebras remains largely unexplored.

One of the main objectives of this paper is to study finite-dimensional Koszul algebras \(\Lambda\) via their Koszul duals \(\Lambda^{!}\). In what follows, we show that this invariant is closely related to a natural homological invariant on the Koszul dual side.

\medskip 
Let \(M=\bigoplus_{i\in\mathbb Z} M_i\) be a finite-dimensional graded
\(\Lambda\)-module. Set
\[
\deg_{\min}(M)=\inf\{\, i\in\mathbb Z \mid M_i\neq 0 \,\},
\qquad
\deg_{\max}(M)=\sup\{\, i\in\mathbb Z \mid M_i\neq 0 \,\}.
\]
The \emph{graded length} of \(M\) is defined by
\[
\ell_{\mathrm{gr}}(M)
=
 \deg_{\max}(M)-\deg_{\min}(M) + 1 .
\]

Let \(I^{\bullet}\) be a bounded colinear complex in
\[
\mathcal{CL}^{b}(\Lambda^{!}).
\]
Suppose that the graded module
\[
\bigoplus_{n\in\mathbb Z} H^{n}(I^{\bullet})\langle n\rangle
\]
is finite-dimensional.

Observe that if \(I^{\bullet}\) is such that
\(\mathfrak{F}(I^{\bullet})\cong M\), then, by Proposition~3.4, for all
\(x\in Q_0\) and \(i,m\in\mathbb Z\), there is a natural isomorphism
\[
H^{-i}(I^{\bullet})_m(x)
\cong
D\!\left(
\operatorname{Ext}^{i-m}_{\Lambda\textup{-}\mathrm{gmod}}
\bigl(M,S_x\langle m\rangle\bigr)
\right).
\]
In particular, for every \(x\in Q_0\) and every \(i\in\mathbb Z\),
\[
H^{-i}(I^{\bullet})_i(x)
\cong
D\!\left(
\operatorname{Hom}_{\Lambda\textup{-}\mathrm{gmod}}
\bigl(M,S_x\langle i\rangle\bigr)
\right).
\]
\begin{lemma}
Let \(M\) be a finite-dimensional graded \(\Lambda\)-module of finite
projective dimension. Then
\[
\operatorname{pd}_{\Lambda}(M)
=
\mathrm{CHL}_{\mathrm{gr}}(I^{\bullet})-1,
\]
where
\[
I^{\bullet}\in
\mathsf{D}^{b}\!\bigl(
\Lambda^{!}\textup{-}\mathrm{Cop}^{\mathbb{Z}}
\bigr)
\]
is such that \(\mathfrak{F}(I^{\bullet})\cong M\).
\end{lemma}

\begin{proof}
Let
\[
P^{\bullet}\longrightarrow M
\]
be the minimal graded projective resolution of \(M\). By Proposition~3.5,
there exists a bounded colinear complex
\[
I^{\bullet}\in
\mathsf{D}^{b}\!\bigl(
\Lambda^{!}\textup{-}\mathrm{Cop}^{\mathbb{Z}}
\bigr)
\]
with finite-dimensional cohomology such that
\[
\mathfrak{F}(I^{\bullet})\cong M
\]
and
\[
\mathfrak{F}\!\left(
\bigoplus_{i\in\mathbb Z}
H^{i}(I^{\bullet})[-i]
\right)
\cong
\operatorname{Lin}^{\Lambda}(P^{\bullet}).
\]

Since \(P^{\bullet}\) is minimal, its linear part has the same nonzero
cohomological range as \(P^{\bullet}\). Under graded Koszul duality, this
range is measured by the graded support of
\[
\bigoplus_{i\in\mathbb Z}
H^{i}(I^{\bullet})\langle i\rangle.
\]
Hence
\[
\operatorname{pd}_{\Lambda}(M)+1
=
\ell_{\mathrm{gr}}\!\left(
\bigoplus_{i\in\mathbb Z}
H^{i}(I^{\bullet})\langle i\rangle
\right)
=
\mathrm{CHL}_{\mathrm{gr}}(I^{\bullet}),
\]
and therefore
\[
\operatorname{pd}_{\Lambda}(M)
=
\mathrm{CHL}_{\mathrm{gr}}(I^{\bullet})-1.
\]
\end{proof}
The following result concerning the graded finitistic dimension of a finite-dimensional Koszul algebra is now an immediate consequence of Lemma~3.9.

\begin{theorem}
Let \(\Lambda\) be a finite-dimensional Koszul algebra with Koszul dual \(\Lambda^{!}\). Then
\[
\mathrm{grfin\,dim}(\Lambda)
=
\sup
\left\{
\mathrm{CHL}_{\mathrm{gr}}(I^{\bullet})
\;\middle|\;
I^{\bullet}\in
\mathcal{CL}^{b}(\Lambda^{!})
\right\}
-1.
\]
In particular, the graded finitistic dimension of \(\Lambda\) is finite if and only if
\[
\sup
\left\{
\mathrm{CHL}_{\mathrm{gr}}(I^{\bullet})
\;\middle|\;
I^{\bullet}\in
\mathcal{CL}^{b}(\Lambda^{!})
\right\}
< \infty.
\]
\end{theorem}
In~\cite{12}, Green et al.\ showed that the graded finitistic dimension is finite for the category of Koszul modules. Since the graded derived Koszul duality \(\mathfrak{F}\) sends coKoszul modules to Koszul modules, while its quasi-inverse \(\mathfrak{G}\) sends Koszul modules to coKoszul modules; see~\cite[Proposition~3.1]{5}, it follows that the graded lengths of coKoszul modules over \(\Lambda^{!}\) are uniformly bounded. We therefore obtain the following result.

\begin{corollary}
Let \(\Lambda\) be an absolutely Koszul algebra such that \(\mathrm{gl\,ld}_{\Lambda}<\infty\). Then
\[
\mathrm{grfin\,dim}(\Lambda)<\infty.
\]
\end{corollary}

\begin{proof}
Assume that \(\Lambda\) is a finite-dimensional absolutely Koszul algebra with \(\mathrm{gl\,ld}_{\Lambda}=n\). Let \(M\) be a finite-dimensional graded \(\Lambda\)-module with finite projective dimension. Then the \(n\)-th syzygy \(\Omega^{n}(M)\) is weakly Koszul.

Let \(I^{\bullet}\) be the bounded colinear complex associated to \(\Omega^{n}(M)\) via the graded derived Koszul duality. By the characterization of weakly Koszul modules, it follows that all cohomology objects \(H^{j}(I^{\bullet})\) are colinear for every \(j\in\mathbb{Z}\).

By~\cite[Theorem~4.5]{12}, the finitistic dimension of Koszul modules is finite. Since the graded derived Koszul duality functor \(\mathfrak{F}\) sends coKoszul modules to minimal projective resolutions of Koszul modules, it follows that coKoszul modules have uniformly bounded graded length. In particular, there exists a uniform bound on the graded lengths of the cohomology objects \(H^{j}(I^{\bullet})\).

Consequently, the graded cohomological length \(\mathrm{CHL}_{\mathrm{gr}}(I^{\bullet})\) is uniformly bounded. The desired conclusion now follows from Theorem~3.10, namely,
\[
\mathrm{grfin\,dim}(\Lambda)<\infty.
\]
\end{proof}
\begin{remark}
By \cite{5}, quadratic monomial algebras are absolutely Koszul and have global linearity defect at most one. In particular, their graded finitistic dimension is finite.
\end{remark}
We now use the characterization of absolutely Koszul algebras in terms of their Koszul duals in order to refine the Koszul dualities established in~\cite{4}. We begin by recalling these dualities in the general case, without any co-coherence assumption on \(\Lambda^{!}\).

\begin{theorem}[{\cite[Theorems~3.8, 3.16, 3.17, 4.3, 4.6, 4.8]{4}}]
Let \(\Lambda\) be a finite-dimensional Koszul algebra with Koszul dual \(\Lambda^{!}\). Then the following hold.

\begin{enumerate}\itemsep=0.4em

\item[\textup{(i)}] \emph{Graded derived Koszul duality.}
There is a triangulated equivalence
\[
\mathfrak{F} \colon
\mathsf{D}^{b}\bigl(\Lambda^{!}\textup{-Cop}^{\mathbb{Z}}\bigr)
\;\xrightarrow{\ \sim\ }\;
\mathsf{D}^{b}\bigl(\Lambda\textup{-gmod}\bigr).
\]

\item[\textup{(ii)}] \emph{Graded singular Koszul duality.}
There is a triangulated equivalence
\[
\mathsf{D}^{b}\!\Bigl(\Lambda^{!}\textup{-Cop}^{\mathbb{Z}}\big/\Lambda^{!}\textup{-gmod}\Bigr)
\;\xrightarrow{\ \sim\ }\;
\mathsf{D}_{\mathrm{sg}}\bigl(\Lambda\textup{-gmod}\bigr).
\]
If, in addition, \(\Lambda\) is Iwanaga--Gorenstein, then there is a triangulated equivalence
\[
\mathsf{D}^{b}\!\Bigl(\Lambda^{!}\textup{-Cop}^{\mathbb{Z}}\big/\Lambda^{!}\textup{-gmod}\Bigr)
\;\xrightarrow{\ \sim\ }\;
\Lambda\textup{-}\underline{\mathrm{Gproj}}^{\mathbb{Z}}.
\]

\item[\textup{(iii)}] \emph{Ungraded derived and singular Koszul dualities.}
There are triangulated equivalences
\[
\mathrm{H}^{0}\!\Bigl(
\operatorname{pretr}\bigl(
\mathsf{D}_{\mathrm{dg}}^{b}(\Lambda^{!}\textup{-Cop}^{\mathbb{Z}})
\big/\langle 1 \rangle[-1]
\bigr)
\Bigr)
\;\xrightarrow{\ \sim\ }\;
\mathsf{D}^{b}\bigl(\Lambda\textup{-mod}\bigr),
\]
\[
\mathrm{H}^{0}\!\Bigl(
\operatorname{pretr}\bigl(
\mathsf{D}_{\mathrm{dg}}^{b}(\Lambda^{!}\textup{-Cop}^{\mathbb{Z}}\big/\Lambda^{!}\textup{-gmod})
\big/\langle 1 \rangle[-1]
\bigr)
\Bigr)
\;\xrightarrow{\ \sim\ }\;
\mathsf{D}_{\mathrm{sg}}\bigl(\Lambda\textup{-mod}\bigr).
\]
If, in addition, \(\Lambda\) is Iwanaga--Gorenstein, then there is a triangulated equivalence
\[
\mathrm{H}^{0}\!\Bigl(
\operatorname{pretr}\bigl(
\mathsf{D}_{\mathrm{dg}}^{b}(\Lambda^{!}\textup{-Cop}^{\mathbb{Z}}\big/\Lambda^{!}\textup{-gmod})
\big/\langle 1 \rangle[-1]
\bigr)
\Bigr)
\;\xrightarrow{\ \sim\ }\;
\Lambda\textup{-}\underline{\mathrm{Gproj}}.
\]

\end{enumerate}

Here \(\mathrm{H}^{0}\!\bigl(\operatorname{pretr}(-)\bigr)\) denotes the homotopy category of the pretriangulated hull of the corresponding dg orbit category; see~\cite[Subsection~2.4]{4} for further details.
\end{theorem}
Now suppose that \(\Lambda\) is absolutely Koszul. Then the category of coperfect graded \(\Lambda^{!}\)-modules coincides with the category of finitely copresented graded \(\Lambda^{!}\)-modules. In this case, the Koszul dualities recalled above admit the following refinement, as announced in the introduction.

\begin{corollary}
Let \(\Lambda\) be a finite-dimensional absolutely Koszul algebra. Then the following triangulated equivalences hold.

\begin{enumerate}\itemsep=0.4em

\item[\textup{(i)}] \emph{Graded derived Koszul duality.}
\[
\mathsf{D}^{b}\bigl(\Lambda^{!}\textup{-Fcp}^{\mathbb{Z}}\bigr)
\;\xrightarrow{\ \sim\ }\;
\mathsf{D}^{b}\bigl(\Lambda\textup{-gmod}\bigr).
\]

\item[\textup{(ii)}] \emph{Graded singular Koszul duality.}
\[
\mathsf{D}^{b}\!\Bigl(
\Lambda^{!}\textup{-Fcp}^{\mathbb{Z}} \big/ \Lambda^{!}\textup{-gmod}
\Bigr)
\;\xrightarrow{\ \sim\ }\;
\mathsf{D}_{\mathrm{sg}}\bigl(\Lambda\textup{-gmod}\bigr).
\]

\item[\textup{(iii)}] \emph{Graded BGG correspondence.}
If \(\Lambda\) is Iwanaga--Gorenstein, then there is a triangulated equivalence
\[
\mathsf{D}^{b}\!\Bigl(
\Lambda^{!}\textup{-Fcp}^{\mathbb{Z}} \big/ \Lambda^{!}\textup{-gmod}
\Bigr)
\;\xrightarrow{\ \sim\ }\;
\Lambda\textup{-}\underline{\mathrm{Gproj}}^{\mathbb{Z}}.
\]

\item[\textup{(iv)}] \emph{Ungraded derived Koszul duality.}
\[
\mathrm{H}^{0}\!\Bigl(
\operatorname{pretr}\bigl(
\mathsf{D}_{\mathrm{dg}}^{b}\bigl(\Lambda^{!}\textup{-Fcp}^{\mathbb{Z}}\bigr)
\big/ \langle 1 \rangle[-1]
\bigr)
\Bigr)
\;\xrightarrow{\ \sim\ }\;
\mathsf{D}^{b}\bigl(\Lambda\textup{-mod}\bigr).
\]

\item[\textup{(v)}] \emph{Ungraded singular Koszul duality.}
\[
\mathrm{H}^{0}\!\Bigl(
\operatorname{pretr}\bigl(
\mathsf{D}_{\mathrm{dg}}^{b}\bigl(
\Lambda^{!}\textup{-Fcp}^{\mathbb{Z}} \big/ \Lambda^{!}\textup{-gmod}
\bigr)
\big/ \langle 1 \rangle[-1]
\bigr)
\Bigr)
\;\xrightarrow{\ \sim\ }\;
\mathsf{D}_{\mathrm{sg}}\bigl(\Lambda\textup{-mod}\bigr).
\]

\item[\textup{(vi)}] \emph{Ungraded BGG correspondence.}
If \(\Lambda\) is Iwanaga--Gorenstein, then there is a triangulated equivalence
\[
\mathrm{H}^{0}\!\Bigl(
\operatorname{pretr}\bigl(
\mathsf{D}_{\mathrm{dg}}^{b}\bigl(
\Lambda^{!}\textup{-Fcp}^{\mathbb{Z}} \big/ \Lambda^{!}\textup{-gmod}
\bigr)
\big/ \langle 1 \rangle[-1]
\bigr)
\Bigr)
\;\xrightarrow{\ \sim\ }\;
\Lambda\textup{-}\underline{\mathrm{Gproj}}.
\]

\end{enumerate}
\end{corollary}
Motivated by the preceding results, we formulate the following conjecture, which is known to hold for quadratic monomial algebras and, more generally, for monomial algebras.

\begin{conjecture}
Let \(\Lambda\) be a finite-dimensional graded algebra. If the graded finitistic dimension of \(\Lambda\) is finite, then the finitistic dimension of \(\Lambda\) is finite.
\end{conjecture}

In the monomial case, syzygies admit a particularly simple description, which allows one to establish the conjecture. For general finite-dimensional Koszul algebras, however, no comparably explicit description of syzygies is available, and the problem remains open.

\medskip 

We conclude the paper by answering a question raised by Green et al.~\cite{12}. Recall that a finite-dimensional graded \(\Lambda\)-module \(M\) is said to have a \emph{linear presentation} if it admits a presentation
\[
P^{-1}\longrightarrow P^{0}\longrightarrow M\longrightarrow 0,
\]
where \(P^{0}\) is generated in degree \(0\) and \(P^{-1}\) is generated in degree \(-1\).

\medskip 

Green et al.~\cite{12} observed that a necessary and sufficient condition on
\(\Lambda\) ensuring that the category of modules with linear presentations,
denoted by \(\mathcal{L}(\Lambda)\), coincides with the category of Koszul
modules, denoted by \(\mathcal{K}(\Lambda)\), was not known. Several necessary
or sufficient conditions were subsequently obtained; see, for
instance,~\cite{1,13}. A complete characterization, however, remained open.
We provide such a characterization in terms of the Koszul dual algebra
\(\Lambda^{!}\).

In fact, the finite-dimensionality hypothesis on \(\Lambda\) may be replaced
by local finite-dimensionality, and the result remains valid in this more
general setting.

\begin{theorem}
Let \(\Lambda\) be a finite-dimensional Koszul algebra. Then the following
conditions are equivalent:
\begin{enumerate}
\item[(a)] \(\mathcal{L}(\Lambda)=\mathcal{K}(\Lambda)\).

\item[(b)] Every graded \(\Lambda^{!}\)-module concentrated in degrees
\(-1\) and \(0\) is weakly coKoszul.

\item[(c)] Every graded \(\Lambda^{!}\)-module concentrated in degrees
\(-1\) and \(0\) can be completed to a coKoszul module.
\end{enumerate}
\end{theorem}

\begin{proof}
We only prove the equivalence between \textup{(a)} and \textup{(b)}, since the proof of the equivalence between \textup{(a)} and \textup{(c)} is analogous.

Assume that
\(\mathcal{L}(\Lambda)=\mathcal{K}(\Lambda)\), and let \(N\) be a graded
\(\Lambda^{!}\)-module concentrated in degrees \(-1\) and \(0\). Under
Koszul duality, \(\mathfrak{F}(N)\) is a two-term linear complex of graded
projective \(\Lambda\)-modules. In particular, it has the form
\[
P^{-1}\longrightarrow P^{0},
\]
and \(H^{0}(\mathfrak{F}(N))\) has a linear presentation. Hence
\[
H^{0}(\mathfrak{F}(N))\in\mathcal{L}(\Lambda)
=\mathcal{K}(\Lambda).
\]
Moreover, since the complex is linear, its other nonzero cohomology module
\(H^{-1}(\mathfrak{F}(N))\) is, up shift, a
syzygy of \(H^{0}(\mathfrak{F}(N))\), and is therefore Koszul. Thus all
cohomology modules of \(\mathfrak{F}(N)\) are Koszul. By the
characterization of weakly coKoszul modules under graded Koszul duality,
\(N\) is weakly coKoszul. This proves \textup{(a)} \(\Rightarrow\)
\textup{(b)}.

Conversely, assume \textup{(b)}, and let
\(M\in\mathcal{L}(\Lambda)\). Choose a linear presentation
\[
P^{-1}\longrightarrow P^{0}\longrightarrow M\longrightarrow0.
\]
The corresponding two-term linear complex is of the form
\(\mathfrak{F}(N)\) for a graded \(\Lambda^{!}\)-module \(N\) concentrated
in degrees \(-1\) and \(0\). By assumption, \(N\) is weakly coKoszul.
Hence the zero cohomology module of \(\mathfrak{F}(N)\) is Koszul. In
particular,
\[
M\cong H^{0}(\mathfrak{F}(N))
\]
is Koszul. Therefore
\(\mathcal{L}(\Lambda)\subseteq\mathcal{K}(\Lambda)\). The reverse
inclusion is immediate, since every Koszul module has a linear
presentation. Consequently,
\[
\mathcal{L}(\Lambda)=\mathcal{K}(\Lambda).
\]
Thus \textup{(a)} and \textup{(b)} are equivalent.
\end{proof}
\begin{thebibliography}{99}

\bibitem{1}
R.~M.~Aquino and E.~L.~Green,
\emph{On modules with linear presentations over Koszul algebras},
Comm. Algebra \textbf{33} (2005), no.~1, 19--36.
\url{https://doi.org/10.1081/AGB-200036765}.

\bibitem{2}
M.~Artin and J.~J.~Zhang,
\emph{Noncommutative projective schemes},
Adv. Math. \textbf{109} (1994), 228--287.

\bibitem{3}
A.~Beilinson, V.~Ginzburg, and W.~Soergel,
\emph{Koszul duality patterns in representation theory},
J. Amer. Math. Soc. \textbf{9} (1996), no.~2, 473--527.

\bibitem{4}
A.~M.~Bouhada,
\emph{A non-graded Koszul duality and its applications},
preprint (2026), available at
\url{https://arxiv.org/abs/2604.16805}.

\bibitem{5}
A.~M.~Bouhada,
\emph{Koszul duality for quadratic monomial algebras},
preprint (2026), available at
\url{https://arxiv.org/abs/2604.20177}.

\bibitem{6}
A.~M.~Bouhada,
\emph{Koszul duality for coherent sheaves},
preprint (2026), available at
\url{https://arxiv.org/abs/2607.14299}.

\bibitem{7}
A.~M.~Bouhada, M.~Huang, and S.~Liu,
\emph{Koszul duality for non-graded derived categories},
preprint (2019), available at
\url{https://arxiv.org/abs/1908.06153}.

\bibitem{8}
A.~M.~Bouhada, M.~Huang, Z.~Lin, and S.~Liu,
\emph{A representation theoretic perspective of Koszul theory},
preprint (2024), available at
\url{https://arxiv.org/abs/2411.15449}.

\bibitem{9}
A.~Conca, E.~De~Negri, and M.~E.~Rossi,
\emph{Koszul algebras and regularity},
in \emph{Commutative Algebra: Expository Papers Dedicated to David Eisenbud on the Occasion of his 65th Birthday},
I.~Peeva (ed.), Springer, 2013, pp.~285--315.

\bibitem{10}
A.~Conca, S.~B.~Iyengar, H.~D.~Nguyen, and T.~R\"omer,
\emph{Absolutely Koszul algebras and the Backelin--Roos property},
Acta Math. Vietnam. \textbf{40} (2015), no.~3, 353--374.

\bibitem{11}
D.~Eisenbud, G.~Fl{\o}ystad, and F.-O.~Schreyer,
\emph{Sheaf cohomology and free resolutions over exterior algebras},
Trans. Amer. Math. Soc. \textbf{355} (2003), 4397--4426.

\bibitem{12}
E.~L.~Green, R.~Mart\'inez-Villa, I.~Reiten, {\O}.~Solberg, and D.~Zacharia,
\emph{On modules with linear presentations},
J. Algebra \textbf{205} (1998), no.~2, 578--604.

\bibitem{13}
E.~L.~Green, E.~Marcos, and P.~Zhang,
\emph{Koszul modules and modules with linear presentations},
Comm. Algebra \textbf{31} (2003), no.~6, 2745--2770.
\url{https://doi.org/10.1081/AGB-120021891}.

\bibitem{14}
J.~Herzog and S.~Iyengar,
\emph{Koszul modules},
J. Pure Appl. Algebra \textbf{201} (2005), no.~1--3, 154--188.

\bibitem{15}
S.~B.~Iyengar and T.~R\"omer,
\emph{Linearity defects of modules over commutative rings},
J. Algebra \textbf{322} (2009), 3212--3237.

\bibitem{16}
D.~Madsen,
\emph{The Koszul dual of a weakly Koszul module},
J. Algebra \textbf{284} (2005), 479--496.

\bibitem{17}
R.~Mart\'inez-Villa and A.~Martsinkovsky,
\emph{Cohomology of tails, Tate--Vogel cohomology, and noncommutative Serre duality over Koszul quiver algebras},
J. Algebra \textbf{280} (2004), 58--83.

\bibitem{18}
R.~Mart\'inez-Villa and M.~Saor\'in,
\emph{Koszul equivalences and dualities},
Pacific J. Math. \textbf{214} (2004), no.~2, 359--378.

\bibitem{19}
R.~Mart\'inez-Villa and D.~Zacharia,
\emph{Approximations with modules having linear resolutions},
J. Algebra \textbf{266} (2003), 671--697.

\bibitem{20}
V.~Mazorchuk, S.~Ovsienko, and C.~Stroppel,
\emph{Quadratic duals, Koszul dual functors, and applications},
Trans. Amer. Math. Soc. \textbf{361} (2009), no.~3, 1129--1172.

\bibitem{21}
I.~Mori,
\emph{Rationality of the Poincar\'e series for Koszul algebras},
J. Algebra \textbf{276} (2004), 602--624.

\bibitem{22}
D.~Piontkovski,
\emph{Coherent algebras and noncommutative projective lines},
J. Algebra \textbf{319} (2008), 3280--3290.

\bibitem{23}
A.~Polishchuk,
\emph{Noncommutative proj and coherent algebras},
Math. Res. Lett. \textbf{12} (2005), no.~1, 63--74.

\bibitem{24}
I.~Reiten and M.~Van den Bergh,
\emph{Noetherian hereditary abelian categories satisfying Serre duality},
J. Amer. Math. Soc. \textbf{15} (2002), no.~2, 295--366.

\bibitem{25}
C.~A.~Weibel,
\emph{An Introduction to Homological Algebra},
Cambridge Studies in Advanced Mathematics, vol.~38,
Cambridge University Press, 1994.

\bibitem{26}
K.~Yanagawa,
\emph{Castelnuovo--Mumford regularity for complexes and weakly Koszul modules},
J. Pure Appl. Algebra \textbf{213} (2009), no.~3, 399--413.

\bibitem{27}
K.~Yanagawa,
\emph{Linearity defect and regularity over a Koszul algebra},
Math. Scand. \textbf{104} (2009), no.~2, 273--290.

\end{thebibliography}
\end{document}
